\begin{frame}{Git là gì? Vì sao dùng Git?}
\begin{itemize}
  \item Git là hệ quản trị phiên bản \textbf{phân tán}, do Linus Torvalds tạo năm 2005.
  \item Mỗi bản sao là một kho đầy đủ lịch sử, làm việc offline rồi đồng bộ sau.
  \item Mọi thay đổi được ghi theo \texttt{commit}, mỗi commit có mã SHA định danh duy nhất.
  \item Gắn với ASPX: quản lý Web Forms, \texttt{Web.config}, biết rõ ai sửa gì, khi nào.
\end{itemize}
\vspace{0.3em}
\begin{tabular}{lll}
\toprule
Tiêu chí & Tập trung (SVN) & Phân tán (Git) \\
\midrule
Nơi lưu lịch sử & Chỉ trên server & Mọi máy đều đủ \\
Mất mạng & Không commit được & Commit bình thường \\
Sập server & Mất lịch sử & Khôi phục từ máy khác \\
\bottomrule
\end{tabular}
\end{frame}

\begin{frame}{Kiến trúc 3 vùng của Git}
\begin{itemize}
  \item \textbf{Working Directory}: các file đang sửa trực tiếp trên đĩa.
  \item \textbf{Staging Area} (Index): vùng đệm chứa các thay đổi chờ ghi nhận.
  \item \textbf{Repository} (thư mục \texttt{.git}): kho lưu toàn bộ lịch sử commit.
  \item \texttt{git add} đưa file từ Working Directory vào Staging Area.
  \item \texttt{git commit} ghi nội dung Staging Area thành một commit trong Repository.
  \item \texttt{git status} cho biết từng file đang nằm ở vùng nào.
  \item Ví dụ ASPX: sửa \texttt{Default.aspx} rồi add và commit để lưu phiên bản.
\end{itemize}
\end{frame}

\begin{frame}[fragile]{Luồng lệnh Git cơ bản}
\begin{itemize}
  \item Cấu hình danh tính một lần duy nhất trên mỗi máy.
  \item Khởi tạo kho mới bằng \texttt{init}, hoặc tải kho có sẵn bằng \texttt{clone}.
  \item Vòng lặp hằng ngày: sửa file, \texttt{add}, \texttt{commit}, xem \texttt{log}.
\end{itemize}
\begin{lstlisting}[language=bash]
git config --global user.name "Nguyen Van A"
git config --global user.email "a@sv.edu.vn"
git init AspxDemo_NET481
cd AspxDemo_NET481
git add Default.aspx Web.config
git commit -m "feat: them trang Default"
git log --oneline
\end{lstlisting}
\end{frame}

\begin{frame}[fragile]{Nhánh và hợp nhất (branch và merge)}
\begin{itemize}
  \item \textbf{Branch} là con trỏ nhẹ tới một commit, dùng để phát triển song song.
  \item Nhánh \texttt{main} luôn ổn định; mỗi tính năng làm trên nhánh \texttt{feature} riêng.
  \item \texttt{merge} gộp nhánh tính năng trở lại nhánh chính sau khi xong việc.
  \item Nhóm 3 người: mỗi TV một nhánh feature, gộp qua Pull Request.
\end{itemize}
\begin{lstlisting}[language=bash]
git switch -c feature/login
# sua file Login.aspx roi commit
git add Login.aspx
git commit -m "feat: them trang dang nhap"
git switch main
git merge feature/login
git branch -d feature/login
\end{lstlisting}
\end{frame}

\begin{frame}{Xung đột (conflict) và cách giải quyết}
\begin{itemize}
  \item Conflict xảy ra khi 2 nhánh cùng sửa một chỗ của một file rồi merge.
  \item Git đánh dấu vùng xung đột bằng \texttt{HEAD}, \texttt{====} và tên nhánh.
\end{itemize}
Quy trình xử lý 4 bước:
\begin{enumerate}
  \item Mở file, đọc kỹ 2 phía và giữ lại nội dung đúng, xóa dấu đánh dấu.
  \item Kiểm tra kỹ, đặc biệt với file \texttt{Web.config} và chuỗi kết nối.
  \item Đánh dấu đã xong bằng lệnh \texttt{git add} trên file vừa sửa.
  \item Hoàn tất bằng \texttt{git commit} (hoặc \texttt{git merge --continue}).
\end{enumerate}
\end{frame}

\begin{frame}[fragile]{Tệp gitignore cho project ASP.NET}
\begin{itemize}
  \item \texttt{.gitignore} liệt kê file và thư mục Git phải bỏ qua.
  \item Project ASPX 4.8.1 phải bỏ qua \texttt{bin}, \texttt{obj}, \texttt{packages} và file user.
  \item Không commit chuỗi kết nối thật và file \texttt{.user} của Visual Studio.
  \item File đã commit rồi thì \texttt{.gitignore} không còn tác dụng với nó.
\end{itemize}
\begin{lstlisting}[language=bash]
# Build output
[Bb]in/
[Oo]bj/
# NuGet packages
packages/
*.nupkg
# Visual Studio user files
*.user
*.suo
.vs/
# Publish output
PublishProfiles/
App_Data/*.mdf
App_Data/*.ldf
\end{lstlisting}
\end{frame}

\begin{frame}{Demo live của TV1: checklist 6 việc}
\begin{itemize}
  \item Demo trên project \texttt{AspxDemo} NET Framework 4.8.1 thật, khoảng 5--7 phút.
\end{itemize}
\begin{enumerate}
  \item \texttt{init} và \texttt{clone} kho demo, cấu hình danh tính và remote.
  \item Sửa \texttt{Default.aspx}, thực hiện \texttt{add}, \texttt{commit}, xem \texttt{log}.
  \item Tạo nhánh \texttt{feature}, commit riêng rồi \texttt{merge} về \texttt{main}.
  \item Tạo conflict cố ý trên \texttt{Web.config} và giải quyết đủ 4 bước.
  \item Demo \texttt{stash}, \texttt{reset} và \texttt{revert} khi sửa nhầm code.
  \item Gắn \texttt{tag} phiên bản v1.0 và đẩy toàn bộ lên GitHub.
\end{enumerate}
\end{frame}
